\chapter{Induction}

Induction is the main tool for proving statements of the form ``for all $n \in \N$, $P(n)$''.
Math 135 also uses it to justify recursive definitions and to prove the binomial theorem.

\section{Sigma and product notation}

$\sum_{i=1}^n a_i = a_1 + a_2 + \cdots + a_n$ and $\prod_{i=1}^n a_i = a_1 a_2 \cdots a_n$.
Standard manipulations: $\sum (a_i + b_i) = \sum a_i + \sum b_i$, $\sum c a_i = c \sum a_i$,
and the index can be shifted: $\sum_{i=1}^n a_i = \sum_{j=0}^{n-1} a_{j+1}$.
An empty sum is $0$ and an empty product is $1$.

\section{The principle of mathematical induction}

\begin{theorem}[Principle of Mathematical Induction (POMI)]
Let $P(n)$ be an open sentence with domain $\N$. If
\begin{enumerate}
  \item (base case) $P(1)$ is true, and
  \item (inductive step) for every $k \in \N$, $P(k) \Rightarrow P(k+1)$,
\end{enumerate}
then $P(n)$ is true for all $n\in\N$.
\end{theorem}

The base case can start at any integer $b$, giving $P(n)$ for all $n \ge b$. In the inductive
step, the assumption $P(k)$ is called the \emph{inductive hypothesis} (IH).

\begin{important}[Template]
\textbf{Base case.} Verify $P(b)$ explicitly (compute both sides).\\
\textbf{Inductive hypothesis.} Assume $P(k)$ holds for some $k \ge b$.\\
\textbf{Inductive step.} Show $P(k+1)$, pointing out where the IH is used.\\
\textbf{Conclusion.} By POMI, $P(n)$ holds for all $n \ge b$.
\end{important}

\begin{example}
Prove $\sum_{i=1}^n i^2 = \dfrac{n(n+1)(2n+1)}{6}$ for all $n \in \N$.

\emph{Base case} $n = 1$: LHS $= 1$, RHS $= \frac{1\cdot 2\cdot 3}{6} = 1$.

\emph{Inductive step.} Assume $\sum_{i=1}^k i^2 = \frac{k(k+1)(2k+1)}{6}$ for some $k\in\N$. Then
\begin{align*}
\sum_{i=1}^{k+1} i^2 &= \frac{k(k+1)(2k+1)}{6} + (k+1)^2 &&\text{(IH)}\\
&= \frac{(k+1)\bigl[k(2k+1) + 6(k+1)\bigr]}{6} = \frac{(k+1)(2k^2 + 7k + 6)}{6}\\
&= \frac{(k+1)(k+2)(2k+3)}{6},
\end{align*}
which is the formula with $n = k+1$. By POMI the formula holds for all $n$.
\end{example}

\begin{example}[Divisibility]
Prove $3 \mid 4^n - 1$ for all $n \ge 0$.
Base: $4^0 - 1 = 0 = 3\cdot 0$. Step: assume $4^k - 1 = 3m$. Then
$4^{k+1} - 1 = 4\cdot 4^k - 1 = 4(3m + 1) - 1 = 12m + 3 = 3(4m+1)$.
\end{example}

\begin{example}[Inequality with a later base case]
Prove $2^n > n^2$ for all $n \ge 5$.
Base: $2^5 = 32 > 25$. Step: assume $2^k > k^2$ with $k \ge 5$. Then
$2^{k+1} = 2\cdot 2^k > 2k^2 = k^2 + k^2$. Since $k \ge 5$, $k^2 \geq 5k = 2k + 3k > 2k + 1$, so
$2^{k+1} > k^2 + 2k + 1 = (k+1)^2$. (The claim fails for $n = 4$: $16 = 16$.)
\end{example}

\section{Strong induction}

\begin{theorem}[Principle of Strong Induction (POSI)]
Let $P(n)$ be an open sentence. If $P(b), P(b+1),\dots,P(b + r)$ are true (several base cases)
and for every $k \ge b + r$,
\[ P(b)\wedge P(b+1)\wedge\cdots\wedge P(k) \Rightarrow P(k+1), \]
then $P(n)$ holds for all $n \ge b$.
\end{theorem}

Strong induction is needed when $P(k+1)$ depends on some earlier case other than $P(k)$, such
as in recurrences of depth $2$ or in factorization arguments. The number of base cases must be
enough that the inductive step only ever refers to cases already covered.

\begin{example}[Recurrence]
Define $x_1 = 3$, $x_2 = 5$, and $x_n = 3x_{n-1} - 2x_{n-2}$ for $n \ge 3$. Prove
$x_n = 2^n + 1$ for all $n \in \N$.

\emph{Base cases:} $x_1 = 3 = 2^1 + 1$ and $x_2 = 5 = 2^2 + 1$.

\emph{Inductive step:} let $k \ge 2$ and assume $x_i = 2^i + 1$ for all $1 \le i \le k$. Then
\begin{align*}
x_{k+1} = 3x_k - 2x_{k-1} &= 3(2^k + 1) - 2(2^{k-1} + 1) \\
&= 3\cdot 2^k - 2^k + 3 - 2 = 2\cdot 2^k + 1 = 2^{k+1}+1.
\end{align*}
Both $k$ and $k-1$ are in $\{1,\dots,k\}$ because $k \ge 2$; this is why two base cases are
needed. Check: $x_3 = 15 - 6 = 9 = 2^3 + 1$.
\end{example}

\begin{example}[Fibonacci bound]
Let $f_1 = f_2 = 1$ and $f_n = f_{n-1} + f_{n-2}$. Prove $f_n < 2^n$.
Base: $f_1 = 1 < 2$, $f_2 = 1 < 4$. Step ($k \ge 2$): $f_{k+1} = f_k + f_{k-1} < 2^k + 2^{k-1}
< 2^k + 2^k = 2^{k+1}$.
\end{example}

\begin{proposition}
Every integer $n \ge 2$ can be written as a product of primes.
\end{proposition}
\begin{proof}
Strong induction on $n$. Base: $2$ is prime. Suppose every integer $m$ with $2 \le m \le k$
is a product of primes, and consider $k+1$. If $k + 1$ is prime we are done. Otherwise
$k + 1 = ab$ with $2 \le a, b \le k$. By the IH both $a$ and $b$ are products of primes, so
$k+1 = ab$ is too.
\end{proof}

\section{The well-ordering principle}

\begin{theorem}[Well-Ordering Principle (WOP)]
Every non-empty subset of $\N$ has a smallest element.
\end{theorem}

WOP is equivalent to induction. It gives an alternative ``minimal counterexample'' style of
proof: if $P(n)$ failed for some $n$, take the \emph{smallest} such $n$ and derive a
contradiction, usually by producing a smaller counterexample.

\begin{example}
Reprove that every $n \geq 2$ is a product of primes with WOP. Suppose not; let $S$ be the
non-empty set of counterexamples and $n = \min S$. Then $n$ is not prime, so $n = ab$ with
$1 < a,b < n$. By minimality, $a,b \notin S$, so both are products of primes, and so is $n$.
Contradiction.
\end{example}

\section{Recursively defined sequences}

A sequence is defined \emph{recursively} by giving initial terms and a rule expressing $x_n$
in terms of earlier terms. A common task is to guess a closed form by computing a few terms
and then prove it by (strong) induction.

\begin{example}
$a_1 = 1$, $a_{n} = 2a_{n-1} + 1$ for $n \geq 2$. Terms: $1, 3, 7, 15, 31$, suggesting
$a_n = 2^n - 1$. Induction: $a_1 = 1 = 2^1 - 1$, and if $a_k = 2^k - 1$ then
$a_{k+1} = 2(2^k - 1) + 1 = 2^{k+1} - 1$.
\end{example}

\section{The binomial theorem}

\begin{definition}
For integers $0 \le k \le n$, $\displaystyle\binom nk = \frac{n!}{k!(n-k)!}$, the number of
$k$-element subsets of an $n$-element set. By convention $\binom nk = 0$ if $k < 0$ or
$k > n$.
\end{definition}

\begin{lemma}[Pascal's identity]
For $n \geq 1$ and $1 \leq k \leq n$: $\binom{n}{k} = \binom{n-1}{k-1} + \binom{n-1}{k}$.
\end{lemma}
\begin{proof}
For $1 \le k \le n-1$,
\[
\binom{n-1}{k-1} + \binom{n-1}{k} = \frac{(n-1)!}{(k-1)!(n-k)!} + \frac{(n-1)!}{k!(n-1-k)!}
= \frac{(n-1)!\,[k + (n-k)]}{k!(n-k)!} = \binom nk .
\]
For $k = n$ both sides equal $1$.
\end{proof}

\begin{theorem}[Binomial theorem]
For all $n \ge 0$ and $a,b$ (real or complex),
\[ (a+b)^n = \sum_{k=0}^n \binom nk a^{n-k} b^k. \]
\end{theorem}
\begin{proof}
Induction on $n$. For $n = 0$ both sides are $1$. Assume the result for $n$. Then
\begin{align*}
(a+b)^{n+1} &= (a+b)\sum_{k=0}^n\binom nk a^{n-k}b^k
 = \sum_{k=0}^n\binom nk a^{n+1-k}b^k + \sum_{k=0}^{n}\binom nk a^{n-k}b^{k+1}\\
&= \sum_{k=0}^{n+1}\binom nk a^{n+1-k}b^k + \sum_{k=0}^{n+1}\binom n{k-1}a^{n+1-k}b^k
 \qquad\text{(shift index; added terms are $0$)}\\
&= \sum_{k=0}^{n+1}\left[\binom nk + \binom n{k-1}\right]a^{n+1-k}b^k
 = \sum_{k=0}^{n+1}\binom{n+1}k a^{n+1-k}b^k
\end{align*}
by Pascal's identity (the $k=0$ term is $\binom n0 = 1 = \binom{n+1}0$).
\end{proof}

\begin{example}
The coefficient of $x^5$ in $(2x - 3)^8$: the general term is $\binom 8k (2x)^{8-k}(-3)^k$;
we need $8 - k = 5$, so $k = 3$, giving $\binom83 2^5 (-3)^3 = 56\cdot 32\cdot(-27) = -48384$.
\end{example}

\begin{example}
Setting $a = b = 1$: $\sum_{k=0}^n \binom nk = 2^n$. Setting $a = 1$, $b = -1$ (for $n\geq 1$):
$\sum_{k=0}^n (-1)^k\binom nk = 0$.
\end{example}

\section{Common pitfalls}
\begin{itemize}
  \item Skipping or fudging the base case. The ``proof'' that all horses are the same colour
        fails because the step $P(1) \Rightarrow P(2)$ is invalid.
  \item Assuming what you need to prove (``assume $P(n)$ for all $n$'').
  \item Using too few base cases in strong induction.
  \item Not stating where the inductive hypothesis is used.
\end{itemize}

\section{Summary}
\begin{itemize}
  \item POMI: base case plus $P(k) \Rightarrow P(k+1)$. POSI: enough base cases plus
        $P(b)\wedge\cdots\wedge P(k) \Rightarrow P(k+1)$.
  \item WOP (every non-empty subset of $\N$ has a least element) is equivalent and supports
        minimal-counterexample proofs.
  \item Recurrences: guess a closed form, prove it by strong induction.
  \item Binomial theorem $(a+b)^n = \sum \binom nk a^{n-k}b^k$, proved via Pascal's identity.
\end{itemize}
